\documentclass[11pt]{article}
\usepackage[margin=1in]{geometry}
\usepackage{amsmath,amssymb,bm,braket}
\usepackage{hyperref}
\begin{document}
\paragraph*{Problem setup:}
In goniopolar materials, the thermal power can be n-type in one direction and p-type in another direction.

\paragraph*{Main problem:}
Consider a two-band 2D intrinsic semiconductor under temperature $T$. The conduction band has a dispersion $E_c= \Delta/2+\sum_{\alpha=x,y}\frac{\hbar^2k_{\alpha}^2}{2m_{c,\alpha}}$, and the valence band has a dispersion $E_v= -\Delta/2-\sum_{\alpha=x,y}\frac{\hbar^2k_{\alpha}^2}{2m_{v,\alpha}}$, where $\Delta$ is the size of band gap. We further assume that the relaxation time of the electrons and holes are the same. What condition should the effective mass $m_{c/v,\alpha}$, with $ \alpha=x,y$, satisfy to exhibit goniopolarity?



\paragraph*{Transport regime}
The equal electron and hole relaxation times already specified above are
energy independent. Use the nondegenerate semiclassical Boltzmann regime,
with intrinsic carrier neutrality and constant relaxation time. Distinguish
the condition valid throughout this regime from a further large-gap
approximation $\Delta\gg k_BT$; if that approximation is used, state it.
\end{document}
